% conclusion.tex
This project set out to build a complete, reproducible fMRI pipeline
in MATLAB and SPM25, and to test it on two structurally different
task paradigms. Looking at the results as a whole, the pipeline
delivered what it was supposed to.

Preprocessing quality was good across both datasets. Motion was low,
coregistration overlays showed accurate alignment, and tSNR values
were in the normal range for 3\,T acquisitions. The GLM analyses
produced spatially sensible activation maps: the three-condition
somatotopic motor map for Dataset~1 matched the known homuncular
organisation of primary motor cortex almost exactly, and the
category-selective visual responses in Dataset~2 — particularly the
face-selective fusiform cluster — replicated a well-known result from
the original Haxby et~al.\ study~\cite{Haxby2001}. Functional
connectivity analyses revealed coherent system-level networks anchored
to the seed regions, and the cross-dataset comparison confirmed that
the two datasets were comparable in data quality while differing
appropriately in their functional topography.

From an engineering perspective, the key design decision was to keep
all processing scripted and sequential (seven numbered MATLAB files),
with skip-if-exists logic so that individual steps can be re-run or
debugged without reprocessing the entire pipeline. Every intermediate
output is saved with a consistent naming scheme, making it easy to
identify which file came from which processing stage.

There is clearly room to extend this work. The analysis here covers
a single subject per dataset; a proper group study would require
second-level random-effects models and family-wise or FDR correction
for multiple comparisons. The connectivity analysis could be made
more rigorous by regressing out global signal and white matter / CSF
noise before computing correlations. Dynamic connectivity approaches
could reveal how network structure evolves across the scan. But for
the scope of this course project, the pipeline achieves its goals:
it works, it is reproducible, and the results make neurobiological
sense.
